\chapter{在群论中的应用}
\label{chapter:ch14}

前一章中我们讨论了如何将拓扑问题 (即空间 $B$ 的所有覆叠空间的分类问题) 化为代数问题 (即空间 $B$ 的基本群的所有子群的分类问题). 现在我们来考虑这个问题的反问题, 即将代数问题化为拓扑问题. 这个代数问题是证明自由群的子群还是自由群. 这个结论好像没有什么问题, 但是却没有什么人给出过简洁的证明. 我们将从应用覆叠空间理论于某些称为线性图的拓扑空间入手来处理它.

在第 83 节中, 我们定义线性图并证明其覆叠空间仍是线性图. 第 84 节证明连通图的基本群是自由群, 自由生成元组一一对应于极大树之外的边. 第 85 节利用覆叠空间证明自由群的子群仍是自由群, 并给出指数公式: 若 $F$ 有 $n+1$ 个自由生成元, 子群 $H$ 的指数为 $k$, 则 $H$ 有 $kn+1$ 个自由生成元.

\subfile{Chapter-14/section-83}
\subfile{Chapter-14/section-84}
\subfile{Chapter-14/section-85}
